\RequirePackage[l2tabu, orthodox]{nag} % warn about outdated packages
\documentclass[12pt,leqno]{article}
\usepackage{microtype} %
\usepackage{setspace}
\onehalfspacing
% Let TeX loosen interword spacing rather than overflow the margin when a long
% unbreakable math atom lands awkwardly.
\emergencystretch=3em
\usepackage{xcolor, color, ucs}     % http://ctan.org/pkg/xcolor
\usepackage{natbib}
\usepackage{booktabs}          % package for thick lines in tables
\usepackage{amsfonts,amsthm,amsmath,amssymb}          % AMS Stuff
\usepackage{empheq}            % To use left brace on {align} environment
\usepackage{graphicx}          % Insert .pdf, .eps or .png
\usepackage{enumitem}          % http://ctan.org/pkg/enumitem
\usepackage[mathscr]{euscript}          % Font for right expectation sign
\usepackage{tabularx}          % Get scale boxes for tables
\usepackage{float}             % Force floats around
\usepackage{afterpage}% http://ctan.org/pkg/afterpage
\usepackage[T1]{fontenc}
\usepackage{rotating}          % Rotate long tables horizontally
\usepackage{bbm}                % for bold betas
\usepackage{bm} 
\usepackage{csquotes}           % \enquote{} and \textquote[][]{} environments
\usepackage{subfig}
\usepackage{lscape}
\usepackage{titling}            % modify maketitle in latex
% \usepackage{mathtools}          % multlined environment with size option
\usepackage{verbatim}
\usepackage{geometry}
\usepackage{bigfoot}
\usepackage[format=hang,
            font={small},
            labelfont=bf,
            textfont=rm]{caption}
\usepackage{tikz}
\usetikzlibrary{positioning}

\geometry{verbose,margin=2cm,nomarginpar}
\setcounter{secnumdepth}{2}
\setcounter{tocdepth}{2}

\usepackage{url}
\usepackage{relsize}            % \mathlarger{} environment

% COLOR FOR LINKS (defined here because hyperref reads it as an option below)
\definecolor{linknavy}{RGB}{31,78,121}

\usepackage[unicode=true,
            pdfusetitle,
            bookmarks=true,
            bookmarksnumbered=true,
            bookmarksopen=true,
            bookmarksopenlevel=2,
            breaklinks=false,
            pdfborder={0 0 1},
            backref=page,
            colorlinks=true,
            hyperfootnotes=true,
            hypertexnames=false,
            pdfstartview={XYZ null null 1},
            citecolor=linknavy,
            linkcolor=linknavy,
            urlcolor=linknavy]{hyperref}
\usepackage{hypernat}

\usepackage{multirow}
\usepackage{titlesec}

\titleformat*{\section}{\large\bfseries}
\titleformat*{\subsection}{\bfseries}
\usepackage[noabbrev]{cleveref} % Should be loaded after \usepackage{hyperref}

\parskip=12pt
\parindent=0pt
\delimitershortfall=-1pt
\interfootnotelinepenalty=100000

\makeatletter
\def\thm@space@setup{\thm@preskip=0pt
\thm@postskip=0pt}
\makeatother

\allowdisplaybreaks
\newtheoremstyle{newstyle}
{12pt} %Aboveskip
{12pt} %Below skip
{\itshape} %Body font e.g.\mdseries,\bfseries,\scshape,\itshape
{} %Indent
{\bfseries} %Head font e.g.\bfseries,\scshape,\itshape
{.} %Punctuation afer theorem header
{ } %Space after theorem header
{} %Heading

\theoremstyle{newstyle}
\newtheorem{thm}{Theorem}
\newtheorem{prop}[thm]{Proposition}
\newtheorem{defin}[thm]{Definition}
\newtheorem{lem}{Lemma}
\newtheorem{cor}{Corollary}
\newcommand{\given}[1][]{\:#1\vert\:}
\DeclareMathOperator{\E}{\mathrm{E}}
\DeclareMathOperator{\Var}{\mathrm{Var}}
\DeclareMathOperator{\Cov}{\mathrm{Cov}}
\newcommand{\R}{\mathbb{R}}              % ordinary symbol, not an operator

% ---------- sources of randomness ----------
\newcommand{\bZ}{\bm{Z}}   \newcommand{\bz}{\bm{z}}   % assignment: random / realized
\newcommand{\bR}{\bm{R}}   \newcommand{\br}{\bm{r}}   % sampling:   random / realized
\newcommand{\EZ}{\E_{\Omega}}    \newcommand{\VarZ}{\Var_{\Omega}}  % over assignment
\newcommand{\ER}{\E_{\Pi}}       \newcommand{\VarR}{\Var_{\Pi}}     % over sampling

% ---------- potential outcomes (function form) ----------
\newcommand{\po}[2]{y_{#1}(#2)}     % \po{i}{1} -> y_i(1); \po{i}{\bZ} -> y_i(Z)
\newcommand{\pov}[1]{\bm{y}(#1)}    % \pov{1}   -> bold y(1)
\newcommand{\ybar}[1]{\bar{y}(#1)}  % \ybar{1}  -> sample mean of y_i(1)

% ---------- estimands and estimator ----------
\newcommand{\tauS}{\tau_{\text{SATE}}}
\newcommand{\tauP}{\tau_{\text{PATE}}}
\newcommand{\tauhat}{\hat{\tau}}

% ---------- variance families: letter = divisor family, subscript = index set ----------
\newcommand{\Ssq}[2][n]{S_{#1}^{2}(#2)}
\newcommand{\ssq}[2][n]{\sigma_{#1}^{2}(#2)}
\newcommand{\scov}[2][n]{\sigma_{#1}(#2)}
\newcommand{\Ssqhat}[2][n]{\widehat{S}_{#1}^{2}(#2)}
\newcommand{\ssqhat}[2][n]{\widehat{\sigma}_{#1}^{2}(#2)}
\newcommand{\Varhat}{\widehat{\Var}}

% ---------- misc ----------
\newcommand{\ind}[1]{\mathbbm{1}\left[#1\right]}


% \Pr is a limits-taking operator by default, so a subscript like \Pr_{\Omega}
% lands fully beneath Pr in display style; redefine it so the distribution
% subscript sits beside the operator, matching \E and \Var.
\renewcommand{\Pr}{\operatorname{Pr}}

% COLORS FOR EQUATIONS (3-class Dark2, colorblind-safe)
\definecolor{eqgreen}{RGB}{27,158,119}
\definecolor{eqblue}{RGB}{117,112,179}
\definecolor{eqred}{RGB}{217,95,2}

\begin{document}

\title{Unbiasedness of the Difference-in-Means Estimator for the SATE and the PATE\thanks{This is a live document that is subject to updating at any time.}}
\author{Thomas Leavitt}
\date{}
\maketitle

This guide proves that the Difference-in-Means estimator is unbiased for the Sample Average Treatment Effect (SATE), under both complete and simple random assignment, and, when the study sample is itself a random draw from a superpopulation, that it is also unbiased for the Population Average Treatment Effect (PATE). It is self-contained, but it shares its framework with the companion guides \enquote{Notation and Setup for Design-Based Causal Inference} (the full development of the notation used here), \enquote{Variance of the Difference-in-Means Estimator for the SATE and Estimation of That Variance,} and \enquote{Variance of the Difference-in-Means Estimator for the PATE and Estimation of That Variance.}

\section{Estimation of Sample Average Treatment Effect (SATE)}

\subsection{Setup} \label{sec: SATE setup}

This section states the framework only as far as the proofs require; the companion guide \enquote{Notation and Setup for Design-Based Causal Inference} develops each object in full, with the reasoning behind it.

Let the index $i \in \{1, \ldots , n\}$ run over the $n$ units of a finite sample, $\mathcal{S}_n$, of which $n_T \geq 2$ receive treatment and $n_C = n - n_T \geq 2$ receive control, so that $n \geq 4$ (these bounds matter only for the variance estimator in the companion guides; the setup guide explains them). The indicator $Z_i \in \{0, 1\}$ records whether unit $i$ is assigned to treatment $(Z_i = 1)$ or control $(Z_i = 0)$, and the random assignment vector $\bZ = \begin{bmatrix} Z_1, \dots , Z_n\end{bmatrix}^\top$ collects the indicators. Every assignment vector determines how many units it sends to treatment: Write $n_T(\bz) \coloneqq \sum_{i = 1}^n z_i$ for this count and $n_C(\bz) \coloneqq n - n_T(\bz)$ for its complement, so that $n_T(\bZ)$ is a random variable. Under \textit{complete random assignment} the design fixes the count: The support of $\bZ$ is $\Omega = \{\bz \in \{0, 1\}^n: n_T(\bz) = n_T\}$, on which $n_T(\bZ)$ equals the design constant $n_T$ with probability one, and $\bZ$ is distributed uniformly on $\Omega$, whose cardinality\footnote{For an arbitrary set $W$, let $\lvert W \rvert$ denote the cardinality of (i.e., the number of elements in) the set $W$.} is $\lvert \Omega \rvert = \binom{n}{n_T}$. Under \textit{simple random assignment} no constraint fixes the count, and $n_T(\bZ)$ is a nondegenerate random variable; \Cref{sec: simple assign} treats that case by conditioning on the event $\{n_T(\bZ) = n_T\}$.

Under the Stable Unit Treatment Value Assumption (SUTVA), unit $i$'s potential outcome depends on the assignment vector only through unit $i$'s own coordinate, so each unit carries two \textit{fixed} potential outcomes, $\po{i}{1}$ and $\po{i}{0}$, with individual effect $\tau_i \coloneqq \po{i}{1} - \po{i}{0}$. The observed outcome is $Y_i \coloneqq \po{i}{Z_i} = Z_i \po{i}{1} + (1 - Z_i)\po{i}{0}$, random only through $Z_i$. The setup guide develops the potential outcomes schedule, states SUTVA precisely, and walks through the collapse from $\po{i}{\bZ}$ to $\po{i}{1}$ and $\po{i}{0}$ in numbered steps.

Write $\ybar{1} \coloneqq n^{-1} \sum_{i = 1}^n \po{i}{1}$ and $\ybar{0} \coloneqq n^{-1} \sum_{i = 1}^n \po{i}{0}$ for the finite population means of the treatment and control potential outcomes; both are fixed quantities. The target of interest is the Sample Average Treatment Effect (SATE),
\begin{equation*}
\tauS \coloneqq n^{-1} \sum \limits_{i = 1}^n \tau_i = n^{-1} \sum \limits_{i = 1}^n \left(\po{i}{1} - \po{i}{0}\right) = \ybar{1} - \ybar{0},
\end{equation*}
which is likewise a fixed (though unknown) quantity. Define the Difference-in-Means estimator of $\tauS$ as
\begin{equation}
\label{eq: diff means est I SATE}
\tauhat \coloneqq \dfrac{\sum \limits_{i = 1}^n Z_i Y_i}{\sum\limits_{i = 1}^n Z_i}  - \dfrac{\sum \limits_{i = 1}^n \left(1 - Z_i\right) Y_i}{\sum\limits_{i = 1}^n \left(1 - Z_i\right)}.
\end{equation}
Unlike the estimand $\tauS$, the estimator $\tauhat$ is a \textit{random variable}, since it is a function of the random assignment vector $\bZ$ and inherits all of its randomness from $\bZ$. The denominators of the ratio form are the counts $n_T(\bZ)$ and $n_C(\bZ)$; under complete random assignment the design fixes them at $n_T$ and $n_C$, so $\tauhat$ can equivalently be written
\begin{equation*}
\tauhat = \dfrac{1}{n_T} \sum \limits_{i = 1}^n Z_i Y_i - \dfrac{1}{n_C} \sum \limits_{i = 1}^n \left(1 - Z_i\right) Y_i,
\end{equation*}
which is the form used in the companion guides on the estimator's variance. For the expectation of $\tauhat$ in \Cref{eq: diff means est I SATE}, I write $\EZ[\cdot]$ to emphasize that the expectation is taken over only the randomness of the assignment process, i.e., over the distribution of $\bZ$ on $\Omega$.

\subsection{Estimation under Complete Random Assignment}

%This proof can apply to both complete and simple random assignment at the individual level. Under simple random assignment, the respective numbers of treatment and control units are random quantities. Yet, as is clear from Fisher's example of the purple flower \citep[see][]{little1989,upton1992}, Fisher nevertheless held the view that one ought to treat $n_T$ and $n_C$ as fixed after conditioning on their realized sizes, i.e., on the realized values of the random counts $n_T(\bZ) \coloneqq \sum_{i = 1}^n \ind{Z_i = 1}$ and $n_C(\bZ) \coloneqq \sum_{i = 1}^n \ind{Z_i = 0}$, where $\ind{\cdot}$ is an indicator function that is equal to $1$ if the argument to the function is true and $0$ if it is false. In this proof, we will embrace Fisher's approach such that the proof below is the same regardless of whether $n_T$ and $n_C$ actually are fixed by the random assignment mechanism or one simply conditions upon the realized values of the random counts $n_T(\bZ)$ and $n_C(\bZ)$.

\begin{lem} \label{lem: exp val Z}
Under complete, uniform random assignment in which $n_T$ out of $n$ total units are assigned to treatment, $\EZ[Z_i] = \frac{n_T}{n}$ for every unit $i \in \{1, \dots , n\}$.
\end{lem}
\begin{proof}

We will complete this proof in two steps: We will show that (1) the proportion of assignments in which unit $i$ is in the treatment condition is $\frac{n_T}{n}$ and (2) under uniform, random assignment, the probability that $Z_i = 1$ is equal to this proportion $\frac{n_T}{n}$.

\textbf{Step 1:} The number of ways to choose a subset of $n_T$ treated units from a fixed population of $n$ units is

\begin{equation} \label{eq: total assignments}
\binom{n}{n_T} = \dfrac{n!}{\left(n - n_T\right)!n_T!} = \dfrac{n!}{n_C!n_T!},
\end{equation}
where $n_C = n - n_T$ is the number of units assigned to the control condition.

Given that an arbitrary unit $i$ is in the treatment condition and only $n_T$ total units can be in the treatment condition, there are $\binom{n - 1}{n_T - 1}$ ways in which $n_T - 1$ other units could be in the treatment condition. Hence, the number of assignments in which unit $i$ is treated and $n_T - 1$ other units are treated is:

\begin{equation} \label{eq: assignments unit i treated}
\binom{n - 1}{n_T - 1} = \dfrac{\left(n - 1\right)!}{\left(\left(n - 1\right) - \left(n_T - 1\right)\right)!\left(n_T - 1\right)!}
\end{equation}

To get the proportion of assignments in which unit $i$ is treated, we need to divide \eqref{eq: assignments unit i treated} by \eqref{eq: total assignments}:

\begin{equation} \label{eq: prop assignments unit i treated I}
\dfrac{\binom{n - 1}{n_T - 1}}{\binom{n}{n_T}} = \dfrac{\left(\dfrac{\left(n - 1\right)!}{\left(\left(n - 1\right) - \left(n_T - 1\right)\right)!\left(n_T - 1\right)!}\right)}{\left(\dfrac{n!}{n_C!n_T!}\right)}
\end{equation}

Now notice that:
\begin{align*}
\left(n - 1\right) - \left(n_T - 1\right) & = n - 1 - n_T + 1 \\
& = n - n_T \\
& = n_C
\end{align*}
We can therefore substitute $n_C$ for $(n - 1) - (n_T - 1)$ in \eqref{eq: prop assignments unit i treated I}, which gives us:

\begin{equation}\label{eq: prop assignments unit i treated II}
\dfrac{\left(\dfrac{\left(n - 1\right)!}{n_C!\left(n_T - 1\right)!}\right)}{\left(\dfrac{n!}{n_C!n_T!}\right)} \\
\end{equation}

Now we can simply manipulate \eqref{eq: prop assignments unit i treated II} and cancel terms until we are left with $\frac{n_T}{n}$:
\begin{align*}
& = \left(\dfrac{\left(n - 1\right)!}{n_C!\left(n_T - 1\right)!}\right)\left(\dfrac{n_C!n_T!}{n!}\right) \\
& = \left(\dfrac{\left(n - 1\right)\left(n - 2\right) \dots 1}{n_C\left(n_C - 1\right) \dots 1 \left(n_T - 1\right) \dots 1}\right) \left(\dfrac{n_C\left(n_C - 1\right) \dots 1 n_T \left(n_T - 1\right) \dots 1}{n\left(n - 1\right) \dots 1}\right) \\[2.5ex]
& = \dfrac{\overbrace{\textcolor{eqblue}{\left(n - 1\right)\left(n - 2\right) \dots 2}}^{\text{(a)}} \, \overbrace{\textcolor{eqred}{n_C\left(n_C - 1\right) \dots 2}}^{\text{(b)}} \, n_T \, \overbrace{\textcolor{eqgreen}{\left(n_T - 1\right) \dots 2}}^{\text{(c)}}}{\underbrace{\textcolor{eqred}{n_C\left(n_C - 1\right) \dots 2}}_{\text{(b)}} \, \underbrace{\textcolor{eqgreen}{\left(n_T - 1\right) \dots 2}}_{\text{(c)}} \, n \, \underbrace{\textcolor{eqblue}{\left(n - 1\right) \dots 2}}_{\text{(a)}}} \\
\end{align*}
Each braced group in the numerator has an identical partner in the denominator: (a) with (a), (b) with (b), and (c) with (c). All three pairs cancel, which leaves us with $\frac{n_T}{n}$. Therefore, exactly the fraction $\frac{n_T}{n}$ of all possible assignments are those in which unit $i$ is in the treatment condition.

\textbf{Step 2:} The probability that unit $i$ is treated is the sum of the probabilities of the assignments in which unit $i$ is treated. Under uniform random assignment, each assignment has probability $\frac{1}{\lvert\Omega\rvert}$. Thus,
\begin{align*}
\left(\dfrac{1}{\left\lvert\Omega\right\rvert}\right) \left(\dfrac{n_T}{n}\right)\left\lvert\Omega\right\rvert & = \left(\dfrac{1}{\left\lvert\Omega\right\rvert}\right) \dfrac{n_T\left\lvert\Omega\right\rvert}{n} \\[1ex]
& = \dfrac{n_T\left\lvert\Omega\right\rvert}{\left\lvert\Omega\right\rvert n} \\[1ex]
& = \dfrac{n_T}{n}  
 \end{align*}

Since $\Pr_{\Omega}(Z_i = 1) = \frac{n_T}{n}$ for every unit $i \in \{1, \dots , n\}$, the expected value of $Z_i \in \{0, 1\}$ is $\EZ[Z_i] = 1(\frac{n_T}{n}) + 0(1 - \frac{n_T}{n}) = \frac{n_T}{n}$.
\end{proof}

\begin{prop} \label{prop: complete ran assign}
Under complete, uniform random assignment, $\EZ[\tauhat] = \tauS$.
\end{prop}
\begin{proof}
First, the linearity of expectations implies that
\begin{align*}
\EZ \left[\tauhat\right] & = \EZ\left[\dfrac{\sum \limits_{i = 1}^n Z_i Y_i}{\sum\limits_{i = 1}^n Z_i}  - \dfrac{\sum \limits_{i = 1}^n \left(1 - Z_i\right) Y_i}{\sum\limits_{i = 1}^n \left(1 - Z_i\right)}\right] \\
& = \EZ\left[\dfrac{\sum \limits_{i = 1}^n Z_i Y_i}{\sum\limits_{i = 1}^n Z_i}\right] - \EZ\left[\dfrac{\sum \limits_{i = 1}^n \left(1 - Z_i\right) Y_i}{\sum\limits_{i = 1}^n \left(1 - Z_i\right)}\right]
\end{align*}
and, since complete random assignment fixes the counts at $n_T$ and $n_C$, 
\begin{align*}
\EZ \left[\tauhat\right] & = \dfrac{1}{n_T} \EZ\left[\sum \limits_{i = 1}^n Z_i Y_{i}\right] - \dfrac{1}{n_C}\EZ\left[\sum \limits_{i = 1}^n \left(1 - Z_i\right) Y_{i}\right]
\end{align*}

Since, under SUTVA, the observed outcomes for treated units are equal to those units' treatment potential outcomes, we can substitute $Z_i \po{i}{1}$ for $Z_i Y_i$. Analogously, we can substitute $(1 - Z_i)\po{i}{0}$ for $(1 - Z_i)Y_{i}$. That is, with $Y_i = Z_i \po{i}{1} + (1 - Z_i)\po{i}{0}$, it follows that
\begin{align*}
Z_i Y_i & = Z_i\left(Z_i \po{i}{1} + (1 - Z_i)\po{i}{0}\right) = Z_i \po{i}{1} \text{ and} \\ 
(1 - Z_i) Y_i & = (1 - Z_i)\left(Z_i \po{i}{1} + (1 - Z_i)\po{i}{0}\right) = (1 - Z_i) \po{i}{0},
\end{align*}
which leaves us with
\begin{align*}
\EZ \left[\tauhat\right] & = \dfrac{1}{n_T} \EZ\left[\sum \limits_{i = 1}^n Z_i \po{i}{1}\right] - \dfrac{1}{n_C}\EZ\left[\sum \limits_{i = 1}^n \left(1 - Z_i\right) \po{i}{0}\right] \\
& = \left(\dfrac{1}{n_T} \right) \EZ \left[ Z_1 \po{1}{1} + \dots + Z_n \po{n}{1}\right] - \left(\dfrac{1}{n_C}\right) \EZ\left[\left(1 - Z_1\right) \po{1}{0} + \dots + \left(1 - Z_n\right) \po{n}{0} \right] \\
& = \left(\dfrac{1}{n_T} \right) \left(\EZ \left[ Z_1 \po{1}{1}\right] + \dots + \EZ \left[Z_n \po{n}{1}\right]\right) \\
& \qquad - \left(\dfrac{1}{n_C}\right) \left(\EZ\left[\left(1 - Z_1\right) \po{1}{0}\right] + \dots + \EZ \left[\left(1 - Z_n\right) \po{n}{0} \right]\right) \\
& = \left(\dfrac{1}{n_T} \right) \left(\po{1}{1}\EZ \left[ Z_1 \right] + \dots + \po{n}{1}\EZ \left[Z_n\right]\right) \\
& \qquad - \left(\dfrac{1}{n_C}\right) \left(\po{1}{0}\EZ\left[\left(1 - Z_1\right) \right] + \dots + \po{n}{0}\EZ \left[\left(1 - Z_n\right) \right]\right)
\end{align*}

By Lemma \ref{lem: exp val Z}, $\EZ[Z_i ] = (\frac{n_T}{n})$ for all $i \in \{1, \ldots  , n\}$, which implies that $\EZ[(1 - Z_i) ] = 1 - (\frac{n_T}{n}) = (\frac{n_C}{n})$ for all $i \in \{1, \ldots  , n\}$. Hence, we can substitute $(\frac{n_T}{n})$ for $\EZ[Z_i ]$ and $(\frac{n_C}{n})$ for $\EZ[1 - Z_i ]$, respectively, which then yields
\begin{align*}
= \left(\dfrac{1}{n_T} \right) \left(\dfrac{n_T}{n}\right) \left(\po{1}{1} + \dots + \po{n}{1} \right) - \left(\dfrac{1}{n_C}\right) \left(\dfrac{n_C}{n}\right) \left(\po{1}{0} + \dots + \po{n}{0} \right) \\
= \left(\dfrac{1}{n} \right) \left(\po{1}{1} + \dots + \po{n}{1} \right) - \left(\dfrac{1}{n}\right) \left(\po{1}{0} + \dots + \po{n}{0} \right) \\
= \dfrac{\left(\po{1}{1} + \dots + \po{n}{1} \right)}{n} - \dfrac{\left(\po{1}{0} + \dots + \po{n}{0} \right)}{n} \\
= \ybar{1} - \ybar{0} \\
= \tauS.
\end{align*}
\end{proof}

\subsection{Estimation under Simple Random Assignment} \label{sec: simple assign}

Under simple random assignment, the number of experimental units, $n$, remains fixed, but the arm counts are genuinely random: No design constraint pins $n_T(\bZ) = \sum_{i = 1}^n Z_i$ to a single value, and $n_C(\bZ) = n - n_T(\bZ)$ inherits its randomness. No new notation is needed. The count is random for the same reason $\tauhat$ is, namely that the count is a function of the random vector $\bZ$, and we mark that dependence by displaying the argument. When we condition on the event $\{n_T(\bZ) = n_T\}$, the plain $n_T$ denotes the realized value, fixed from that point on, exactly as $z_i$ denotes the realized value of $Z_i$.

We take the support of $n_T(\bZ)$ to be $\{1, \dots , n - 1\}$, so that neither arm can be empty (an empty arm would leave the estimator $\tauhat$ undefined). The set of possible assignments is therefore $\Omega = \{\bz \in \{0, 1\}^n: 0 < n_T(\bz) < n \}$, which contains $2^n - 2$ elements; the distribution of $\bZ$ is that of $n$ independent Bernoulli draws conditioned on the event $\{0 < n_T(\bZ) < n\}$. Whenever the proof below conditions on the event $\{n_T(\bZ) = n_T\}$, the resulting conditional distribution is uniform on $\{\bz: n_T(\bz) = n_T\}$, which is exactly the complete random assignment distribution; we write the conditioning explicitly rather than changing the symbol $\Omega$.

In the proof that follows, we will draw upon the Law of Iterated Expectations: For two random variables $X$ and $W$ defined on the same probability space, $\E[X] = \E[\E[X \given W]]$. In our finite setting the conditioning variable is the count $n_T(\bZ)$, and the law takes the form of the weighted sum in \Cref{eq: iter expec} below. 

\begin{prop}
\label{prop: simple random assign}
Under simple, uniform random assignment, $\EZ [\tauhat] = \tauS$.
\end{prop}

\begin{proof}
By the law of iterated expectations, the expected value of the Difference-in-Means estimator, $\tauhat$, can be decomposed by conditioning on the realized number of treated units, $n_T(\bZ)$:
\begin{equation}
\label{eq: iter expec}
\begin{split}
\EZ\left[\tauhat\right] & = \EZ\left[\tauhat \given n_T(\bZ) = 1\right] \Pr_{\Omega}\left(n_T(\bZ) = 1\right) \\
& \qquad + \dots + \EZ\left[\tauhat \given n_T(\bZ) = n - 1\right]\Pr_{\Omega}\left(n_T(\bZ) = n - 1\right).
\end{split}
\end{equation}

Conditional on the event $\{n_T(\bZ) = n_T\}$, simple random assignment reduces to complete random assignment with $n_T$ treated units, so Proposition \ref{prop: complete ran assign} applies. By that proposition, the expected value of the estimator conditional on any realized number of treated units $n_T \in \{1, \dots, n - 1\}$ is equal to $\tauS$. Hence, we can rewrite Equation \eqref{eq: iter expec} as
\begin{align*}
\EZ\left[\tauhat\right] & = \tauS\Pr_{\Omega}\left(n_T(\bZ) = 1\right) + \dots + \tauS\Pr_{\Omega}\left(n_T(\bZ) = n - 1\right),
\end{align*}
which we can rewrite as
\begin{align*}
\EZ\left[\tauhat\right] & = \tauS\big[\Pr_{\Omega}\left(n_T(\bZ) = 1\right) + \dots + \Pr_{\Omega}\left(n_T(\bZ) = n - 1\right)\big].
\end{align*}
Finally, note that, since $n_T(\bZ)$ takes values in $\{1, \dots, n - 1\}$ with probability one, the second and third axioms of probability imply that $\big[\Pr_{\Omega}(n_T(\bZ) = 1) + \dots + \Pr_{\Omega}(n_T(\bZ) = n - 1)\big] = 1$. Hence,
\begin{align*}
\EZ\left[\tauhat\right] & = \tauS\big[1\big] \\ \EZ\left[\tauhat\right] & = \tauS,
\end{align*}
which proves the proposition.
\end{proof}

\section{Estimation of the Population Average Treatment Effect (PATE)}

\subsection{Setup}

So far we have treated the $n$ units of $\mathcal{S}_n$ as the entire population of interest and shown that $\tauhat$ is unbiased for the SATE. Suppose now that these $n$ units are themselves a random sample from a larger superpopulation, and ask whether $\tauhat$ is unbiased for the average treatment effect in that superpopulation. (The companion guides develop the full superpopulation framework and the variance of $\tauhat$ about the PATE; here we introduce only what unbiasedness requires.)

Consider a superpopulation $\mathcal{P}_N$ of size $N \geq n$, indexed by $i \in \{1, \ldots, N\}$, in which every unit has fixed potential outcomes $\po{i}{1}$ and $\po{i}{0}$ and a fixed individual effect $\tau_i = \po{i}{1} - \po{i}{0}$. Let $n$ of these $N$ units be drawn into the experimental sample $\mathcal{S}_n$ by simple random sampling, recorded by the indicator $R_i \in \{0, 1\}$, where $R_i = 1$ if unit $i$ is sampled and $R_i = 0$ otherwise; the sampling vector $\bR = \begin{bmatrix} R_1, \dots, R_N \end{bmatrix}^\top$ is distributed uniformly on $\Pi = \{\br \in \{0, 1\}^N : \sum_{i = 1}^N r_i = n\}$. Among the $n$ sampled units, $n_T$ are assigned to treatment and $n_C$ to control by complete random assignment, exactly as before. The estimand is the Population Average Treatment Effect,
\begin{equation*}
\tauP \coloneqq \dfrac{1}{N} \sum \limits_{i = 1}^N \tau_i,
\end{equation*}
a fixed (though unknown) feature of the superpopulation, and the Difference-in-Means estimator becomes
\begin{equation}
\label{eq: diff means est PATE unb}
\tauhat \coloneqq \dfrac{1}{n_T} \sum \limits_{i = 1}^N R_i Z_i Y_i - \dfrac{1}{n_C} \sum \limits_{i = 1}^N R_i \left(1 - Z_i\right) Y_i,
\end{equation}
in which $R_i Z_i = 1$ only for units that are both sampled and treated, $R_i(1 - Z_i) = 1$ only for units that are both sampled and control, and unsampled units ($R_i = 0$) contribute nothing. There are now \textit{two} sources of randomness, the sampling indicators $\{R_i\}_{i = 1}^N$ and the assignment indicators $\{Z_i\}_{i = 1}^n$. We write $\E[\cdot]$ for the expectation over both, $\ER[\cdot]$ for the expectation over the sampling process, and $\EZ[\cdot]$ for the expectation over the assignment process conditional on the realized sample.

\begin{prop} \label{prop: pate unbiased}
Under simple random sampling of $n$ units from $\mathcal{P}_N$ and complete random assignment among the sampled units, $\E[\tauhat] = \tauP$.
\end{prop}

\begin{proof}
We take the expectation over both stages of randomness by iterating, first over the assignment (conditional on the realized sample) and then over the sample,
\begin{align*}
\E\left[\tauhat\right] & = \ER\left[\EZ\left[\tauhat\right]\right].
\end{align*}
Conditional on any realized sample, the $n$ sampled units undergo complete random assignment, so Proposition \ref{prop: complete ran assign} applies \textit{within} the sample and gives $\EZ[\tauhat] = \tauS(\bR)$, where $\tauS(\bR) \coloneqq \frac{1}{n}\sum_{i = 1}^N R_i \tau_i$ is the average effect among the sampled units. Within a fixed sample, as in the earlier sections, this average is a constant; here it is a random variable, because the sampling vector $\bR$ determines which units enter the average, and we display the argument to mark that dependence. Substituting,
\begin{align*}
\E\left[\tauhat\right] & = \ER\left[\tauS(\bR)\right] = \ER\left[\dfrac{1}{n} \sum \limits_{i = 1}^N R_i \tau_i\right].
\end{align*}
Because the individual effects $\tau_i$ are fixed and only the sampling indicators $R_i$ are random, the linearity of expectations gives
\begin{align*}
\E\left[\tauhat\right] & = \dfrac{1}{n} \sum \limits_{i = 1}^N \tau_i \, \ER\left[R_i\right].
\end{align*}
It remains to find $\ER[R_i]$. Under simple random sampling of $n$ units from $N$, the event $\{R_i = 1\}$ plays exactly the role that $\{Z_i = 1\}$ played under complete random assignment of $n_T$ units from $n$. The proportion of samples that include unit $i$ is $\frac{n}{N}$, and under uniform sampling the probability of inclusion equals that proportion. The argument of Lemma \ref{lem: exp val Z}, with $N$ and $n$ in place of $n$ and $n_T$, therefore gives $\ER[R_i] = \frac{n}{N}$ for every $i \in \{1, \ldots, N\}$. Substituting,
\begin{align*}
\E\left[\tauhat\right] & = \dfrac{1}{n} \sum \limits_{i = 1}^N \tau_i \left(\dfrac{n}{N}\right) = \dfrac{1}{N} \sum \limits_{i = 1}^N \tau_i = \tauP,
\end{align*}
which proves the proposition.
\end{proof}

This argument reuses the earlier results rather than repeating them. The complete random assignment proposition handles the inner expectation, and the counting argument of Lemma \ref{lem: exp val Z} handles the outer expectation, now applied to sampling $n$ of $N$ units instead of assigning $n_T$ of $n$. The result also makes precise the sense in which the Difference-in-Means estimator answers two questions at once. With respect to the assignment alone, the estimator is unbiased for the SATE, the average effect among the units actually studied. Once we also average over which units are sampled, the estimator is unbiased for the PATE, the average effect in the superpopulation.

\begin{singlespace}
\bibliographystyle{chicago}
\bibliography{references}
\end{singlespace}

\end{document}